% U9 WS1 -- entropy: sign prediction and absolute entropies. Blocks 1-2.
\documentclass{apchem-worksheet}

\apcunit{9}
\apcunittitle{Thermodynamics and Electrochemistry}
\apcdoctitle{Worksheet 1 \textbullet\ Entropy}
\apcsubtitle{CED 9.1--9.2 \textbullet\ count the gas moles before anything else}

\begin{document}
\wsheader[$\Delta S^\circ = \sum nS^\circ_{\text{prod}} -
\sum nS^\circ_{\text{react}}$ \textbullet\ $S^\circ$ in
\si{\joule\per\kelvin\per\mole}: \ce{H2} 131 \textbullet\ \ce{O2} 205
\textbullet\ \ce{N2} 192 \textbullet\ \ce{H2O(l)} 70 \textbullet\
\ce{H2O(g)} 189 \textbullet\ \ce{NH3} 193 \textbullet\ \ce{NO2} 240
\textbullet\ \ce{N2O4} 304 \textbullet\ \ce{CaCO3} 93 \textbullet\ \ce{CaO}
40 \textbullet\ \ce{CO2} 214.]

\problem[]{\ced{9.1} What entropy measures.}
\begin{parts}
  \item Complete: entropy measures how \answerblank[3.2cm]{dispersed} the
        energy and matter of a system are.
  \item Explain why ``a messy room has high entropy'' is a poor analogy.
        \answerlines[3]{Entropy counts the microscopic arrangements
        available to particles and their energy. Tidying a room rearranges
        macroscopic objects without changing the number of molecular
        arrangements at all, so the analogy describes appearance rather
        than the quantity being measured.}
  \item Rank by absolute entropy at the same temperature: \ce{H2O(s)},
        \ce{H2O(l)}, \ce{H2O(g)}.
        \answerblank[7cm]{\ce{H2O(s)} $<$ \ce{H2O(l)} $<$ \ce{H2O(g)}}
\end{parts}

\problem[]{\ced{9.1} Predict the sign of $\Delta S^\circ$ and give the
controlling reason in three words or fewer:}
\begin{parts}
  \item \ce{2H2(g) + O2(g) -> 2H2O(l)}
        \answerblank[6.4cm]{$-$; 3 mol gas to 0}
  \item \ce{NH4NO3(s) -> NH4+(aq) + NO3^-(aq)}
        \answerblank[6.4cm]{$+$; crystal disperses}
  \item \ce{CO2(s) -> CO2(g)} \answerblank[6.4cm]{$+$; sublimation}
  \item \ce{2SO2(g) + O2(g) -> 2SO3(g)}
        \answerblank[6.4cm]{$-$; 3 mol gas to 2}
  \item \ce{H2(g) + Cl2(g) -> 2HCl(g)}
        \answerblank[6.4cm]{$\approx 0$; 2 mol gas to 2}
\end{parts}

\problem[]{\ced{9.1} Part (e) above has no change in the number of gas
moles. Explain what that implies about the magnitude of $\Delta S^\circ$
and why the sign is hard to predict by inspection.
\answerlines[3]{With the gas count unchanged the dominant contribution
vanishes, so $\Delta S^\circ$ is small --- a few
\si{\joule\per\kelvin} either way. What remains are secondary differences
in molecular complexity and mass, which inspection cannot reliably rank, so
the value must be looked up rather than predicted.}}

\problem[]{\ced{9.1} The second law.}
\begin{parts}
  \item State it as an inequality.
        \answerblank[6.8cm]{$\Delta S_{\text{univ}} > 0$ for a favored
        process}
  \item Water freezing at \SI{-10}{\celsius} is favored, yet the water's
        entropy falls. Resolve this. \answerlines[3]{Freezing releases heat
        to the surroundings. Below \SI{0}{\celsius} the entropy gained by
        the surroundings exceeds the entropy lost by the water, so
        $\Delta S_{\text{univ}}$ remains positive.}
  \item True or false: a reaction with $\Delta S_{\text{sys}} < 0$ can
        never be favored. \answerblank[5.4cm]{false --- see part (b)}
\end{parts}

\problem[]{\ced{9.2} The third law.}
\begin{parts}
  \item A perfect crystal at \SI{0}{\kelvin} has entropy exactly
        \answerblank[1.4cm]{zero}.
  \item $S^\circ$ for \ce{O2(g)} is \SI{205}{\joule\per\kelvin\per\mole},
        but $\Delta H_f^\circ$ for \ce{O2(g)} is zero. Explain why the two
        conventions differ. \answerlines[3]{The third law gives entropy a
        true zero, so absolute entropies can be measured --- and an element
        has no reason to sit at that zero. Enthalpy has no natural origin,
        so only changes are measurable, and elements in their standard
        states are \emph{assigned} zero by convention.}
\end{parts}

\problem[]{\ced{9.2} Calculate $\Delta S^\circ$ for each:}
\begin{parts}
  \item \ce{CaCO3(s) -> CaO(s) + CO2(g)} \workspace[1.8cm]
        \keyonly{\begin{apcanswer}$(40 + 214) - 93 = \mathbf{+161}$~J/K
        \end{apcanswer}}
  \item \ce{N2(g) + 3H2(g) -> 2NH3(g)} \workspace[2cm]
        \keyonly{\begin{apcanswer}$2(193) - [192 + 3(131)] = 386 - 585 =
        \mathbf{-199}$~J/K\end{apcanswer}}
  \item \ce{H2O(l) -> H2O(g)} \answerblank[3.4cm]{$+119$ J/K}
  \item \ce{2NO2(g) -> N2O4(g)} \answerblank[3.4cm]{$-176$ J/K}
  \item \ce{2H2(g) + O2(g) -> 2H2O(l)} \workspace[1.8cm]
        \keyonly{\begin{apcanswer}$2(70) - [2(131) + 205] = 140 - 467 =
        \mathbf{-327}$~J/K\end{apcanswer}}
\end{parts}

\problem[]{\ced{9.2} A student computes $\Delta S^\circ$ for the Haber
process as $193 - (192 + 131) = -130$~J/K. Identify both errors and give
the correct value.
\answerlines[3]{They omitted the coefficient 2 on \ce{NH3} and the
coefficient 3 on \ce{H2}. Correctly
$2(193) - [192 + 3(131)] = \mathbf{-199}$~J/K.}}

\problem[]{\ced{9.2} Check each sign you predicted in problem 2 against the
value you calculated in problem 6 where both exist. State what this
confirms about the gas-mole rule.
\answerlines[3]{Every calculated sign matches the prediction made from the
change in gas moles alone. Where the gas count changes, that term dominates
and the sign can be read off without any table --- which is what makes the
no-calculator multiple-choice questions answerable.}}

\begin{spiralstrip}{Unit 8 \textbullet\ acids and bases}
  \item pH of \SI{0.010}{\Molar} \ce{HCl}: \answerblank[2cm]{2.00}
  \item Buffer method step 1: \answerblank[7.3cm]{stoichiometry to
        completion}
  \item pH $=$ p$K_a$ at which point of a titration?
        \answerblank[3.6cm]{the halfway point}
  \item $K_a \times K_b$ for a conjugate pair: \answerblank[1.6cm]{$K_w$}
  \item Equivalence pH, weak acid $+$ strong base:
        \answerblank[2.2cm]{above 7}
\end{spiralstrip}

\end{document}
